\documentclass[11pt]{article}
\usepackage{preamble}
\setlength{\parskip}{4pt}

\begin{document}

\daybanner{AP Statistics \textperiodcentered\ Unit 1 \textperiodcentered\ Topic 1.12 \textperiodcentered\ B: 2026-10-01 \textperiodcentered\ E: 2026-10-09 \textperiodcentered\ TEACHER KEY}{Two Surveys of Park Fees}

\sectionbanner{CED TETHER}

{\small
\begin{itemize}[leftmargin=1.5em,itemsep=2pt,topsep=2pt]
  \item \textbf{Skill 1.C} --- Describe an appropriate method for gathering and representing data.
  \item \textbf{EU DAT-2} --- The way we collect data influences what we can and cannot say about a population.
  \item \textbf{LO DAT-2.E} --- Identify potential sources of bias in sampling methods. [Skill 1.C]
  \item \textbf{EK DAT-2.E.1} --- Bias occurs when certain responses are systematically favored over others.
  \item \textbf{EK DAT-2.E.2} --- When a sample is comprised entirely of volunteers or people who choose to participate, the sample will typically not be representative of the population (voluntary response bias).
  \item \textbf{EK DAT-2.E.3} --- When part of the population has a reduced chance of being included in the sample, the sample will typically not be representative of the population (undercoverage bias).
  \item \textbf{EK DAT-2.E.4} --- Individuals chosen for the sample for whom data cannot be obtained (or who refuse to respond) may differ from those for whom data can be obtained (nonresponse bias).
  \item \textbf{EK DAT-2.E.5} --- Problems in the data gathering instrument or process result in response bias. Examples include questions that are confusing or leading (question wording bias) and self-reported responses.
  \item \textbf{EK DAT-2.E.6} --- Non-random sampling methods (for example, samples chosen by convenience or voluntary response) introduce potential for bias because they do not use chance to select the individuals.
\end{itemize}
}
\textbf{Essential question:} How do response behavior and question wording limit the estimate a survey can support?\par
\textbf{DOK-3 skill:} 4.A \textperiodcentered\ \textbf{FRQ pattern:} {\small\texttt{name-\allowbreak{}the-\allowbreak{}bias-\allowbreak{}direction-\allowbreak{}and-\allowbreak{}fix}}\par
\textbf{DOK rationale:} Part (c) weighs competing accuracy claims against two different bias mechanisms and defends a targeted repair without assuming an unknown true proportion.\par

\frameworkphaseheader{Do Now (first take) $\rightarrow$ Explore (video + follow-along) $\rightarrow$ Exit (finish + turn in)}{1 $\rightarrow$ 3}{45}{%
\begin{itemize}[leftmargin=1.5em,itemsep=2pt,topsep=2pt]
  \item Board slide up at the bell. Read both council claims without confirming either.
  \item Begin the video follow-along at minute 5.
  \item Return to the board slide at minute 33. Circulate and require a separate diagnosis for each survey.
  \item Collect at the door. Score part (c) only, E/P/I.
\end{itemize}
}{%
\begin{itemize}[leftmargin=1.5em,itemsep=2pt,topsep=2pt]
  \item Choose one estimate or neither and cite one survey detail.
  \item Complete the video follow-along about undercoverage, nonresponse, voluntary response, and response bias.
  \item Finish (a)--(c), revise the first take in the exit line, and turn in the sheet.
\end{itemize}
}{%
\begin{itemize}[leftmargin=1.5em,itemsep=2pt,topsep=2pt]
  \item Who is missing from the mail results even though the original sample was random?
  \item Can you know whether the nonresponders oppose the fee more or less?
  \item Which words in the phone question could move an answer, and in what direction?
  \item Does knowing one bias's likely direction tell you which final estimate is closer?
\end{itemize}
}{Solo teacher. Circulate ~5 pairs. Require students to connect each flaw to a mechanism and a repair; do not accept a percentage choice based only on which sample is larger.}

\begin{callout}{\IconWarn\ WATCH FOR}{calloutred}
\begin{itemize}[leftmargin=1.5em,itemsep=2pt,topsep=2pt]
  \item Confusing nonresponse with voluntary response; the households were selected but did not reply.
  \item Assuming an 80 percent response rate protects against loaded wording.
  \item Claiming the exact size of either bias or ranking closeness without truth data.
\end{itemize}
\end{callout}

\sectionbanner{SCORING PART (c) --- E / P / I}

\textbf{Expected elements}
\begin{itemize}[leftmargin=1.5em,itemsep=2pt,topsep=2pt]
  \item \textbf{rejects-ranking} (required) --- Explains why neither estimate is established as closer to the unknown truth
  \item \textbf{uses-both-designs} (required) --- Links low mail response to possible differences and phone wording to pushed answers
  \item \textbf{one-targeted-fix} (required) --- Gives one workable repair and explains which survey problem it addresses
\end{itemize}

{\small\renewcommand{\arraystretch}{1.25}
\noindent\begin{tabular}{|p{0.5in}|p{5.9in}|}
\hline
\textbf{E} & All three elements are correct and linked to this problem. Accept equivalent plain-language reasoning. \\ \hline
\textbf{P} & At least one element includes a correct evidence-to-reason link, but another is missing or incorrect. \\ \hline
\textbf{I} & No correct evidence-to-reason link; only a choice, copied numbers, or unsupported statements. \\ \hline
\end{tabular}}

\textbf{Common mistakes}
\begin{itemize}[leftmargin=1.5em,itemsep=2pt,topsep=2pt]
  \item Choosing $62\%$ solely because 240 people answered
  \item Saying the mail survey is unbiased because the original 500 were randomly chosen
  \item Assuming mail nonresponse must increase opposition even though its direction is unknown
  \item Fixing the phone survey by calling more numbers but keeping the loaded wording
\end{itemize}

\clearpage
\focusbanner{TODAY'S PROBLEM --- annotated}

Suppose a town studies support for a park fee.\par\smallskip \textbf{Mail survey:} A neutral question is mailed to 500 randomly chosen households. Only 120 return it; $78\%$ of respondents oppose the fee.\par \textbf{Phone survey:} The town calls 300 random phone numbers. Of 240 people who answer, $62\%$ oppose after hearing, \emph{Do you support a NEW FEE that would tax park users?}\par\smallskip One council member says $62\%$ must be closer to the truth because more people answered. Another says $78\%$ must be closer because the households were randomly chosen.\par
\begin{center}
\textbf{Park-fee surveys}\par\smallskip
\begin{tabular}{|l|c|c|c|}
\hline
\textbf{} & \textbf{Selected} & \textbf{Replied} & \textbf{Oppose} \\ \hline
\textbf{Mail} & 500 & 120 & $78\%$ \\ \hline
\textbf{Phone} & 300 & 240 & $62\%$ \\ \hline
\end{tabular}
\end{center}
\begin{firsttakebox}{FIRST TAKE (5 min)}
Which estimate, if either, would you report to the town? Name one detail behind your choice.\par
\answer{Not graded. Either number may feel safer; the goal is to separate known flaws from unknowable bias size.}
\end{firsttakebox}

\textit{Rules callout ``BIAS FAVORS SOME RESPONSES (keep on the page)'' is printed on the student sheet.}\par\medskip

\dokbadge{1}~\textbf{(a)} Find each survey's response rate.\par
\answer{Mail: $120/500=24\%$. Phone: $240/300=80\%$.}

\dokbadge{2}~\textbf{(b)} In two or three sentences, explain one possible bias in each survey. Say which direction each could push opposition, or why its direction is unknown.\par
\answer{Mail nonrespondents may hold systematically different opinions from respondents, so $78\%$ could be too high or too low; the direction is not known from the data. Calling the fee NEW and saying it would tax users likely pushes people toward opposition, so $62\%$ is likely too high as an estimate of opposition to neutrally described policy.}

\dokbadge{3}~\textbf{(c)} Has either council member shown which estimate is closer to the town's true opinion? Use one detail from each survey to justify your answer. Choose one survey to improve, and explain how your change addresses its problem.\par
\answer{Neither claim is established: we do not know the true proportion or the sizes of the biases. Mail respondents may differ from the $76\%$ who did not reply, while the phone wording may push answers toward opposition even with an $80\%$ response rate. Accept either targeted repair: follow up neutrally with selected mail nonrespondents, or replace the phone wording with a neutral question. A repair improves the relevant problem; it does not prove the resulting estimate exact.}

\begin{summaryexitbox}{EXIT}
One thing the video changed about my first take: \hrulefill
\end{summaryexitbox}

\end{document}
